% Reset dei contatori per l'appendice
\setcounter{figure}{0}
\renewcommand{\thefigure}{A.\arabic{figure}}
\setcounter{table}{0}
\renewcommand{\thetable}{A.\arabic{table}}
\section{Ottimizzazione Iperparametrica e Validazione Finale}
\label{app:hyperparam_tuning}

Sulla base dei risultati del benchmark preliminare, i classificatori più performanti sono stati sottoposti a una procedura di ottimizzazione fine Hyperparameter Tuning volta a massimizzare la capacità di generalizzazione e prevenire l'overfitting.

\subsection{Strategia di Ottimizzazione (Randomized Search)}

L'esplorazione dello spazio degli iperparametri \(\Theta\) è stata condotta mediante \textbf{Randomized Search CV}. A differenza della Grid Search esaustiva, questo metodo campiona un numero fisso di configurazioni (\(N_{iter}=230\)) da distribuzioni di probabilità definite a priori. Tale approccio consente di esplorare efficientemente spazi ad alta dimensionalità con un budget computazionale limitato.

Per ogni modello candidato \(M_i\), è stata definita una griglia di ricerca specifica. Nel caso del vincitore \textit{XGBoost}, i parametri critici ottimizzati includono:
\begin{itemize}
    \item \texttt{learning\_rate}: Tasso di apprendimento (shrinkage).
    \item \texttt{max\_depth}: Profondità massima dell'albero (controllo overfitting locale).
    \item \texttt{colsample\_bytree}: Frazione di feature campionate per albero (decorrelazione).
    \item \texttt{subsample}: Frazione di dati campionati (Bagging stocastico).
\end{itemize}

\subsection{Selezione del Modello "Champion"}

La Tabella \ref{tab:final_ranking} riporta la classifica definitiva aggregata, confrontando le varianti Base e Tuned per i modelli al vertice.

\begin{table}[h]
\centering
\begin{tabular}{llcc}
\toprule
\textbf{Rank} & \textbf{Modello} & \textbf{Configurazione} & \textbf{CV F1-Score} \\
\midrule
\textbf{1} & \textbf{XGBoost} & \textbf{Tuned} & \textbf{0.7310} \\
2 & LightGBM & Tuned & 0.7194 \\
3 & XGBoost & Base & 0.7156 \\
4 & LightGBM & Base & 0.7156 \\
5 & Linear SVC & Tuned & 0.6819 \\
6 & AdaBoost & Tuned & 0.6614 \\
\bottomrule
\end{tabular}
\caption{Leaderboard finale post-ottimizzazione. XGBoost Tuned emerge come Best Performer.}
\label{tab:final_ranking}
\end{table}

\noindent
Il modello \textbf{XGBoost (Tuned)} è stato proclamato vincitore con un F1-Score medio di validazione pari a 0.7310. L'ottimizzazione ha prodotto un incremento evidente rispetto alla baseline (+1.54\%), confermando l'efficacia della ricerca degli iperparametri nell'affinare le performance e la stabilità del classificatore.

\subsection{Validazione su Test Set (Hold-out)}

La valutazione finale dell'architettura selezionata è stata eseguita sul Test Set \(S_{test}\) (\(N=6.028\)), mantenuto segregato durante l'intera pipeline di training e tuning. I risultati rappresentano una stima non distorta delle performance attese in produzione.

\begin{itemize}
    \item \textbf{Accuracy:} \(0.8540\) (Il modello classifica correttamente l'85.4\% degli individui).
    \item \textbf{F1-Score (Classe \(>50K\)):} \(0.7330\).
\end{itemize}

\subsection{Analisi degli Errori}
Il report di classificazione evidenzia un comportamento tipico nella gestione di classi non perfettamente bilanciate:
\begin{itemize}
    \item \textbf{Trade-off Precision/Recall:} Il modello dimostra di saper bilanciare in modo ottimale l'affidabilità delle sue previsioni positive (Precision) con la sua capacità di catturare la maggior parte degli individui target (Recall).
\end{itemize}

\noindent
Tale trade-off è perfettamente in linea con l'obiettivo di massimizzare l'F1-Score globale. Il raggiungimento di quota 0.7330 sul Test Set certifica che il modello mantiene un eccellente equilibrio tra la purezza delle predizioni e la copertura del target, confermando la solidità del tuning effettuato.
